В первой задаче требуется по изображению отпечатка пальца предсказать его класс по классификации Генри. Для обучения и оценки используется несколько датасетов:
\begin{itemize}
\item NIST Special Database 4~\cite{NIST-SD-4} - основной датасет. Используется как базовый источник примеров отпечатков для обучения модели предсказанию классов по Генри и для корректного сравнения с работами, опирающимися на стандартизованные наборы NIST.
\item SokoFING~\cite{SOCOFing} - дополнительный датасет. Его полезно использовать в двух режимах: (a) как дополнительный источник обучения (совместно с NIST SD4); (b) как внешний тест, чтобы проверить переносимость модели за пределы основного датасета.
\item Fingerprint Minutiae Computer Vision Dataset~\cite{fingerprint-minutiae-dataset} - датасет c картами минуций. Используется для построения модели, которая будет применяться на этапе предобработки для построения карты направленных минуций. Основная модель будет проводить классификацию уже с использованием карты минуций. 
\item CASIA-FingerprintV5~\cite{CASIA-fingerprintV5} - дополнительный большой датасет для претрейна в self-supervised режиме.
\end{itemize}


Во второй задаче требуется предсказать определенные свойства пористых материалов, в частности, металл-органических каркасных структур (MOF). Чаще всего рассматриваются адсорбционные свойства: например, способность данного MOF поглощать определенный газ при разных давлениях. Ключевой особенностью будет использование топологических дескрипторов, извлечённых из структуры порового пространства материала, в качестве признаков для модели. Для обучения и оценки моделей планируется использовать несколько наборов данных MOF:

\begin{itemize}
\item Hypothetical MOFs~\cite{Wilmer2012} – база гипотетических структур MOF. Этот датасет содержит структуры (54k материалов) из MOFXDB и рассчитанные характеристики адсорбции $\text{CO}_2$ при пяти различных давлениях.
\item Boyd–Woo MOF~\cite{Boyd-Woo} – база предсказанных структур MOF. В этот набор входят рассчитанные показатели адсорбции метана и $\text{CO}_2$ при низком и высоком давлениях, а также коэффициенты Генри для этих газов.
\item CoRE MOF 2019~\cite{CoreMOF2019} – подборка экспериментально синтезированных структур MOF из базы CoRE MOF. Она представляет реальные пористые материалы, пригодные для оценки обобщающей способности моделей.
\end{itemize}


В третьей задаче рассматривается семантическая сегментация медицинских изображений - автоматическое разделение изображения на области, соответствующие различным анатомическим структурам. Акцент делается на топологически корректной сегментации, то есть предсказанные области должны иметь количество связных компонентов и отверстий, совпадающее с анатомически правильным . Для этого планируется встроить сравнение топологических характеристик в лосс при обучении модели. Для обучения и оценки моделей планируется использовать датасет:

\begin{itemize}
\item AMOS~\cite{ji2022amos,amos2022zenodo} - клинический датасет для 3D-сегментации органов брюшной полости на КТ и МРТ. Он включает 600 исследований на 600 пациентах (500 КТ и 100 МРТ), содержит воксельную разметку 15 органов (включая селезёнку, печень, почки, поджелудочную железу, аорту и др.).
\end{itemize}
